%% PM Research Watch - WEEKLY ROLLUP, week ending Saturday 1 August 2026.
%% Authored from templates/weekly.tex. Aggregation only; no API was re-queried.
\documentclass[10pt]{article}
\providecommand{\FIGDIR}{fig}
\providecommand{\DIGESTDATE}{1 August 2026}
\providecommand{\PERIODSTART}{27 July 2026}
\providecommand{\PERIODEND}{1 August 2026}
% ===================== PM Research Watch - shared preamble =====================
\usepackage[T1]{fontenc}
\usepackage[utf8]{inputenc}
\usepackage{amsmath,amssymb}
\usepackage{tgpagella}
\usepackage{tgheros}
\renewcommand{\sfdefault}{qhv}
\renewcommand{\ttdefault}{lmtt}
\usepackage{microtype}
\usepackage[english]{babel}

\usepackage[a4paper,top=17mm,bottom=17mm,left=17mm,right=17mm,
            headsep=6mm,footskip=10mm]{geometry}
\usepackage{graphicx}
\usepackage{xcolor}
\usepackage{colortbl}
\usepackage{tikz}
\usetikzlibrary{calc}
\usepackage[most]{tcolorbox}
\usepackage{titlesec}
\usepackage{fancyhdr}
\usepackage{booktabs}
\usepackage{array}
\usepackage{longtable}
\usepackage{enumitem}
\usepackage{multicol}
\usepackage{ragged2e}
\usepackage[hidelinks]{hyperref}

% ------------------------------------------------------------------ palette
\definecolor{Ink}    {HTML}{14262E}
\definecolor{Deep}   {HTML}{0F4C5C}
\definecolor{Teal}   {HTML}{1F7A8C}
\definecolor{Sky}    {HTML}{4FA3C4}
\definecolor{Sage}   {HTML}{6FA287}
\definecolor{Amber}  {HTML}{E3A008}
\definecolor{Clay}   {HTML}{B5643C}
\definecolor{Coral}  {HTML}{D1495B}
\definecolor{Violet} {HTML}{6A5B9E}
\definecolor{Slate}  {HTML}{4A6572}
\definecolor{Paper}  {HTML}{FBF9F5}
\definecolor{Mist}   {HTML}{EEEAE2}
\definecolor{Rule}   {HTML}{D9D3C9}

\pagecolor{Paper}
\color{Ink}

\hypersetup{colorlinks=true, linkcolor=Deep, urlcolor=Teal, citecolor=Teal,
            pdftitle={PM Research Watch}, pdfauthor={Automated daily digest}}

% ------------------------------------------------------------------ headings
\titleformat{\section}
  {\sffamily\bfseries\large\color{Deep}}{}{0pt}
  {}[{\vspace{1.2mm}\color{Amber}\titlerule[1.1pt]}]
\titlespacing*{\section}{0pt}{5.5mm}{2.8mm}

\titleformat{\subsection}
  {\sffamily\bfseries\normalsize\color{Teal}}{}{0pt}{}
\titlespacing*{\subsection}{0pt}{3.6mm}{1.4mm}

\setlist[itemize]{leftmargin=4.2mm, itemsep=1.1mm, topsep=1.3mm, parsep=0pt,
                  label=\textcolor{Teal}{\footnotesize$\blacktriangleright$}}

% ------------------------------------------------------------------ header/footer
\pagestyle{fancy}
\fancyhf{}
\renewcommand{\headrulewidth}{0pt}
\renewcommand{\footrulewidth}{0pt}
\fancyhead[L]{\sffamily\scriptsize\color{Slate}PM RESEARCH WATCH\;\textcolor{Rule}{\textbar}\;\DIGESTDATE}
\fancyhead[R]{\sffamily\scriptsize\color{Slate}Particulate matter monitoring \& health effects}
\fancyfoot[C]{\sffamily\scriptsize\color{Slate}\thepage}

% ------------------------------------------------------------------ components
% metric tile
\newcommand{\metric}[3]{%
  \begin{tcolorbox}[enhanced, colback=white, colframe=Rule, boxrule=0.5pt,
    arc=1.6mm, left=2.2mm, right=2.2mm, top=1.8mm, bottom=1.8mm,
    borderline west={1.7mm}{0pt}{#3}, nobeforeafter, valign=center,
    height=17mm, width=\linewidth]
    \RaggedRight
    {\sffamily\bfseries\fontsize{17}{17}\selectfont\color{#3}#1}\\[0.5mm]
    {\sffamily\fontsize{6.5}{7.4}\selectfont\color{Slate}#2}
  \end{tcolorbox}}

% pull-quote / key insight box
\newtcolorbox{keybox}[1][]{enhanced, colback=Mist, colframe=Mist,
  arc=1.6mm, left=3.2mm, right=3.2mm, top=2.4mm, bottom=2.4mm,
  borderline west={1.7mm}{0pt}{Amber}, #1}

% subtopic band
\newcommand{\band}[2]{%
  \vspace{3.2mm}\noindent
  \begin{tcolorbox}[enhanced, colback=#2, colframe=#2, arc=1.2mm,
    left=2.6mm, right=2.6mm, top=1.1mm, bottom=1.1mm, boxrule=0pt,
    width=\linewidth, nobeforeafter]
    {\sffamily\bfseries\small\color{white}#1}
  \end{tcolorbox}\vspace{1.6mm}}

% one paper entry: {short cite}{title}{journal}{doi}{body}
\newcommand{\paperentry}[5]{%
  \par\vspace{1.9mm}\noindent
  {\sffamily\bfseries\footnotesize\color{Ink}#1}\;%
  {\sffamily\fontsize{7}{8}\selectfont\color{Slate}#3}\\[0.4mm]
  {\footnotesize\itshape\color{Deep}#2}\\[0.6mm]
  {\footnotesize #4}\\[0.3mm]
  {\fontsize{6.6}{7.6}\selectfont\color{Slate}\texttt{doi:}\,\href{https://doi.org/#5}{\texttt{\detokenize{#5}}}}
  \par}

% caption style
\newcommand{\figcap}[1]{\par\vspace{1.0mm}
  {\fontsize{7.2}{8.6}\selectfont\color{Slate}\textsf{\textbf{Fig.}}\; #1\par}\vspace{1.5mm}}

\providecommand{\permil}{\textperthousand}

\setlength{\parindent}{0pt}
\setlength{\parskip}{0pt}
\renewcommand{\arraystretch}{1.18}


\begin{document}

% ==================================================================== MASTHEAD
\thispagestyle{empty}
\begin{tikzpicture}[remember picture, overlay]
  \fill[Deep] (current page.north west) rectangle
      ($(current page.north east)+(0,-67mm)$);
  \fill[Teal] ($(current page.north west)+(0,-67mm)$) rectangle
      ($(current page.north east)+(0,-68.4mm)$);
  \foreach \x/\y/\r in {22/12/0.9, 41/26/1.5, 63/9/0.7, 88/31/1.1, 112/17/1.9,
                        138/29/0.8, 159/13/1.3, 178/54/1.0, 196/21/1.6,
                        30/44/0.6, 74/57/1.2, 125/48/0.9, 168/6/0.8, 191/59/0.7,
                        52/18/1.0, 100/60/0.6, 148/20/1.1, 66/49/1.4, 205/45/0.8}
    \fill[Sky, opacity=0.30] ($(current page.north west)+(\x mm,-\y mm)$) circle (\r mm);
\end{tikzpicture}

\vspace*{4mm}
{\sffamily\fontsize{7.6}{9}\selectfont\color{Sky}%
 WEEKLY ROLLUP\;\textbullet\;PhD RESEARCH WATCH\par}
\vspace{2.2mm}
{\sffamily\bfseries\fontsize{27}{29}\selectfont\color{white}%
 Particulate Matter\par}
{\sffamily\fontsize{27}{29}\selectfont\color{Sky}%
 Week in Review\par}
\vspace{5.0mm}
{\sffamily\fontsize{8.4}{10}\selectfont\color{white}%
 Week ending Saturday \DIGESTDATE\;\;\textcolor{Teal}{\textbar}\;\;pooled
 \PERIODSTART{} to \PERIODEND\;\;%
 \textcolor{Teal}{\textbar}\;\;131 records from 6 daily issues\par}
\vspace{18mm}

% -------------------------------------------------------------- metric strip
\noindent
\begin{minipage}[t]{0.190\linewidth}\metric{131}{RECORDS POOLED\\THIS WEEK}{Deep}\end{minipage}\hfill
\begin{minipage}[t]{0.190\linewidth}\metric{20}{SENSING RECORDS\\(15\% OF WEEK)}{Teal}\end{minipage}\hfill
\begin{minipage}[t]{0.190\linewidth}\metric{10}{SUBTOPICS\\ACTIVE}{Sky}\end{minipage}\hfill
\begin{minipage}[t]{0.190\linewidth}\metric{48}{POOLED EFFECT\\ESTIMATES}{Amber}\end{minipage}\hfill
\begin{minipage}[t]{0.190\linewidth}\metric{4}{TRIALS ON\\THE WATCH LIST}{Clay}\end{minipage}

\vspace{6mm}

% -------------------------------------------------------------- week in one box
\section*{\sffamily\bfseries\color{Deep}The week in one paragraph}
\vspace{-1.5mm}{\color{Teal}\rule{\linewidth}{1.1pt}}\vspace{2.5mm}

\begin{keybox}
{\sffamily\bfseries\small\color{Deep}Four independent papers this week attacked the
exposure \emph{variable} rather than the health model --- and all four found it
wanting. The through-line is that PM$_{2.5}$ mass at a fixed monitor is now the
weakest link in the causal chain, and the field has started saying so with
numbers.}\\[1.4mm]
\footnotesize
Nyoni et al.\ (31 Jul, \textit{Air Qual Atmos Health}) validated road-geometry proxies
against a year of filter-based personal PM$_{2.5}$ in 343 postpartum participants across
three sub-Saharan countries: weighted road density correlated only in Mozambique
($\rho = 0.351$) and \textbf{distance-to-highway inverted sign in Kenya}
($\rho = -0.224$, $p = 0.018$). Saeedi et al.\ (1 Aug, \textit{ES\&T}) showed
land-use regression for mobile ultrafine particles falling from
\textbf{$\sim$217\% average percentage error under random cross-validation to
$\sim$79\% under spatiotemporal CV} --- the model was never predicting, the validation
was. Blanco et al.\ (1 Aug, \textit{JESEE}) resampled on-road campaign designs against a
5{,}283-person cohort and found every alternative design attenuated the health estimate.
And Zhu et al.\ (31 Jul) reported pneumonia-admission RRs of \textbf{1.004
(0.991--1.017) for PM$_{2.5}$ against 1.221 (1.174--1.271) for NO$_2$} --- PM mass was
the null term in its own model.

What would have to be true for this to be wrong: the four results would have to be
selection rather than signal. Two of them (Saeedi, Blanco) are methods papers whose
premise is that current practice is inadequate, so they are not independent of the
conclusion; Nyoni's setting is one where household biomass dominates the personal
exposure budget, which is exactly where traffic geometry should fail; and Zhu's is a
single-pollutant model over correlated combustion tracers, so it is a statement about
correlation structure, not about which pollutant acts. \textbf{This is a real pattern,
but it is a pattern about measurement adequacy, not about PM being harmless.} The
distinction matters and the week's own papers do not always draw it.
\end{keybox}

\vspace{3mm}

% -------------------------------------------------------------- trend & inflection
\section*{\sffamily\bfseries\color{Deep}Subtopic trend and inflection}
\vspace{-1.5mm}{\color{Teal}\rule{\linewidth}{1.1pt}}\vspace{2.5mm}

\noindent\includegraphics[width=\linewidth]{\FIGDIR/w1_subtopic_trend.png}
\figcap{Per-issue subtopic counts across the week from \texttt{state/metrics.csv}.
\textbf{The one real inflection is sensing/instrumentation: 2, 2, 6, 1, 0, 9 across
27 July--1 August, against daily totals of 33, 19, 25, 18, 11, 25.} Both ends of that
series are pipeline artefacts, not literature. The 6 on 29 July came from a Consensus
backfill; the 0 on 31 July and the 9 on 1 August are the same defect and its fix --- the
Crossref by-journal leg was activated on 1 August and back-swept the week, so nine
records that were deposited between 26 July and 1 August all landed in one issue.
\textbf{Read the weekly total (20/131, 15\%) rather than any daily value.} Every other
subtopic is flat within Poisson noise on a base of 11--33 records/day; nothing else in
this figure is a trend.}

\footnotesize
\textbf{Emerging this week.} Three genuinely new lines. (i) \textbf{Source-resolved
exposure} --- Dales et al.\ separate wildfire-sourced from background PM$_{2.5}$ and find
$+6.07\%$ (5.75--6.39) headache ED visits for the former against a null for the latter;
Zhou M et al.\ attribute 27\% of India's population-weighted PM$_{2.5}$ to transboundary
transport, more than any domestic sector. (ii) \textbf{Ultrafine and number-based
metrics} --- 6 of 131 records this week use UFP or particle-number exposure (Ramsden,
Zhou 2026b, Rusconi, Zhou L, Saeedi, Blanco), against essentially none in the preceding
month's dailies. (iii) \textbf{Oxidative potential as a modelled quantity} --- Mishra
et al.'s KM-OP predicts OP from composition at $R^2 > 0.75$, which is the step that makes
speciated sensing actionable.

\textbf{Fading.} Nothing has genuinely faded in six days, and claiming otherwise would be
noise-fitting. The apparent drop in cardiovascular/metabolic records (4 on 29 July, 2 on
30 July, 0 on 31 July and 1 August) is journal-issue timing on an 11--25-record base, not
a decline.

\textbf{Flat but load-bearing.} Exposure assessment \& modelling (17 records) and
mechanistic toxicology (11) held steady all week and carry the field regardless of the
weekly headline. Reproductive/developmental (10) is the cluster with the most consistent
effect estimates --- Xie's pubertal-timing HRs, Li's stillbirth HRs, Sieber's LBW OR and
Pai's incense aORs are all in the same direction with intervals that exclude the null.

\vspace{3mm}
\noindent\includegraphics[width=\linewidth]{\FIGDIR/f2a_architecture.png}
\figcap{Pooled study architecture. Measurement campaigns lead at 23, but \textbf{reviews
and syntheses are effectively tied at 22 --- one record in six of this week's corpus is a
review}, a high share that reflects a mid-year cluster of review issues rather than new
primary evidence. Observational designs together contribute 36 (16 cohort, 12 acute,
6 cross-sectional, 2 ecological), and \textbf{exactly one record in 131 is a trial or
intervention} --- the evidence base this week is almost entirely observational and
synthetic.}

\clearpage

% -------------------------------------------------------------- pooled analytics
\section*{\sffamily\bfseries\color{Deep}Pooled analytics}
\vspace{-1.5mm}{\color{Teal}\rule{\linewidth}{1.1pt}}\vspace{3mm}

\noindent\includegraphics[width=\linewidth]{\FIGDIR/f1_subtopics.png}
\figcap{Pooled subtopic distribution, \PERIODSTART{} to \PERIODEND. Sensing (20) and
exposure assessment (17) together account for \textbf{28\% of the week} --- the
methodological half of the field is, for once, better represented than any single health
cluster. Neuro/mental health, respiratory/allergic, burden/policy and occupational/indoor
tie at 13.}

\vspace{2mm}
\noindent\includegraphics[width=\linewidth]{\FIGDIR/f2b_endpoint.png}
\figcap{Pooled health endpoint by organ system. \textbf{Thirty-two records carry no human
health endpoint} and eight more were surfaced without an abstract --- together 31\% of the
week, which is what a corpus with a healthy instrumentation and source-attribution
component looks like. Among records that do carry an endpoint, respiratory (17) and
cardiovascular (12) lead; the long tail below 3 records is a single study each and should
not be read as a trend.}

\clearpage

\noindent\includegraphics[width=\linewidth]{\FIGDIR/f5_forest.png}
\figcap{\textbf{Pooled effect estimates for the week --- read as a caterpillar plot, not a
meta-analysis.} No summary diamond is drawn and none should be. \textbf{Pooling is not
defensible here and no $I^2$ is reported}, for four separate reasons: exposure contrasts
differ (per 1\,$\mu$g\,m$^{-3}$ in Xie, per IQR in S\'anchez, a $\geq$9\,$\mu$g\,m$^{-3}$
threshold in Yang, a percentage-change conversion in Dales); outcomes differ across ten
organ systems; several estimates are multiple contrasts from one paper and one cohort
(Yang contributes 3, Xie 3, Ghana AIRCLARE 6, Luo 4), so they are not independent; and
7 of the 48 estimates have \textbf{no published interval at all} (Mokoena $\times$3,
Ouyang $\times$2, Abi Nader $\times$2) and are plotted as points. The one subset that is
internally comparable is the four respiratory RRs per 10\,$\mu$g\,m$^{-3}$ from Zhu and
Yue --- and within that subset the particulate terms (1.004, 1.004) are null while the
gas terms (1.221, 1.068, 1.172) are not.}

\vspace{2mm}
\noindent\includegraphics[width=\linewidth]{\FIGDIR/f4_heatmap.png}
\figcap{Pooled subtopic $\times$ study-architecture cross-tabulation. \textbf{The
persistently empty cells are the finding.} Mechanistic toxicology has no cell in any
observational column --- eleven records, no bridge to human biomarker work. The sensing
row has no cell in any cohort or case-crossover column: across 131 records and six days,
\textbf{no study characterised an instrument and used it to estimate a health effect in
the same paper.} Eight records come closest (Wang's SARS-CoV-2 surveillance, Patton's
outdoor risk assessment, S\'anchez's asthma admissions, Chen's 29-hazard mixture,
Blanco's UFP--cognition analysis) but in every case the exposure instrument is inherited,
not evaluated.}

\clearpage

% -------------------------------------------------------------- groups & citations
\section*{\sffamily\bfseries\color{Deep}Groups, venues and most-cited work}
\vspace{-1.5mm}{\color{Teal}\rule{\linewidth}{1.1pt}}\vspace{2.5mm}

\footnotesize
\textbf{Recurring author and group clusters.} \textbf{There are none, and the apparent
ones are an artefact worth naming.} Surname frequency over the pooled corpus returns
Chen (5), Zhou (5), Ding, Park, Lu, Wang, Zhang and Li (3 each) --- but these are
distinct research groups sharing common surnames, not repeat output. The 1 August issue
alone carries two unrelated first-author Zhous (Liyuan Zhou on Bakhoor emissions, Mi Zhou
on India source attribution), which is why they are disambiguated as ``Zhou L'' and
``Zhou M'' in the register. \textbf{First-author surname is not a group identifier and
should not be used as one}; a real cluster analysis needs ORCID or affiliation strings,
which this pipeline does not currently store. That is a state-schema gap, logged as such.
The only genuine programmatic stream visible is the Henan Rural Cohort (Zhang, 1 Aug) and
the Chongqing Child Health Cohort (Xie, 27 Jul), each contributing one record.

\vspace{2mm}
\textbf{Venue concentration.} Nineteen journals contributed more than one record.
\textit{Toxics} leads at 7, \textit{J Environ Sci (China)}, \textit{J Hazard Mater} and
\textit{Atmos Environ} at 6 each, \textit{Front Public Health} and
\textit{Environ Sci Pollut Res} at 3, \textit{Atmos Chem Phys} at 3. Two of these are
structural rather than substantive: the \textit{Atmos Environ} count is entirely the
Crossref by-journal leg activated on 1 August and reflects a deposit window, not a
research surge; the \textit{Toxics} cluster is a themed-issue effect. \textbf{No
single-journal cluster this week represents an independent replication.}

\vspace{2mm}
\textbf{Most-cited work of the period.} \textbf{Not reported, because it would be noise.}
Every record in this corpus entered an index within the last six days; citation counts
where available are 0--5 and are dominated by how long a preprint sat before formal
publication, not by influence. Ranking a week-old corpus by citations measures
publication lag. The metric becomes meaningful in the monthly and yearly rollups and is
deferred to them.

\vspace{4mm}

% -------------------------------------------------------------- trial watch
\section*{\sffamily\bfseries\color{Deep}Trial registry watch}
\vspace{-1.5mm}{\color{Teal}\rule{\linewidth}{1.1pt}}\vspace{2.5mm}

{\fontsize{7.4}{9}\selectfont\color{Slate}From \texttt{state/trials.json}: registrations
created or updated inside the window. Daily issues accumulate; the rollup is where they
surface.\par}

\vspace{1.5mm}
\footnotesize
\textbf{New this week --- three sham-controlled filtration trials} (surfaced 1 August on
an intervention sweep; all three are indoor HEPA filtration against a sham unit, which is
the correct comparator and rules out the placebo-and-noise objection that sinks
open-label purifier studies).

\begin{itemize}\setlength{\itemsep}{1.0mm}
\item \textbf{NCT05867381} --- \textit{Slowing Atherothrombosis Progression Through Indoor
Air Filtration}, University of Southern California. Recruiting, N/A phase, \textbf{n =
112}, crossover, HEPA vs sham, adults with ischaemic heart disease history, primary
completion \textbf{June 2028}. Crossover on a progression endpoint is the right design
for a within-person exposure contrast, but 112 participants against an atherothrombosis
imaging endpoint is thin: the indoor PM$_{2.5}$ reduction a HEPA unit delivers
($\sim$30--60\% of indoor mass) implies an effect far smaller than the between-person
variance in plaque progression. Expect a surrogate-endpoint result.

\item \textbf{NCT05718245} --- \textit{Using Indoor Air Filtration to Reduce PM$_{2.5}$
Cardiometabolic Effects in At-risk Individuals}, USC. Recruiting, \textbf{n = 52},
HEPA vs sham, primary completion \textbf{December 2026} --- the nearest-term readout on
the list. At $n = 52$ this is powered for biomarker change (blood pressure, glucose,
inflammatory markers), not for events, and should be read that way.

\item \textbf{NCT05016271} --- \textit{Health Benefits of Air Purifiers in Primary School
Students}, Fudan University. Active, not recruiting, \textbf{n = 110}, HEPA-equipped vs
non-HEPA ventilation in classrooms, cardiopulmonary function. \textbf{Primary completion
was June 2022 and no results have been posted} --- four years is well past the reporting
expectation and is itself worth logging.
\end{itemize}

\vspace{1.5mm}
\textbf{Is the endpoint capable of detecting the effect?} For all three: not for hard
outcomes. A sham-controlled filtration trial changes personal PM$_{2.5}$ by single-digit
$\mu$g\,m$^{-3}$ over weeks to months; the cohort literature in this same week
(Luo: CVD-mortality RR 1.014 per 10\,$\mu$g\,m$^{-3}$; Yang: OHCA RR 1.02) implies effect
sizes that a 52--112-person trial cannot resolve on events. \textbf{These trials are
mechanism studies wearing RCT clothing}, and their value is the biomarker and personal-
exposure data, not the clinical endpoint.

\vspace{2mm}
\textbf{Standing watch list.} \textbf{NCT05874479} --- \textit{Reducing Air Pollution
Exposure to Lower Blood Pressure Among NYC Public Housing Residents} (NYU Langone,
$n = 440$, active vs sham portable air cleaner, primary completion September 2027).
Logged 29 July; at $n = 440$ it is the only trial on the list plausibly powered for a
blood-pressure endpoint. Nearest readout to watch: \textbf{NCT05718245, December 2026}.

\clearpage

% -------------------------------------------------------------- gaps & direction
\section*{\sffamily\bfseries\color{Deep}Open gaps}
\vspace{-1.5mm}{\color{Teal}\rule{\linewidth}{1.1pt}}\vspace{2.5mm}

\footnotesize
\begin{itemize}\setlength{\itemsep}{1.4mm}
\item \textbf{Measurement--epidemiology coupling: still zero.} Across 131 records, no
study both characterised an instrument and estimated a health effect with it. Blanco
et al.\ came closest --- a mobile UFP campaign feeding a cognition analysis in the same
paper --- but the monitoring instrument itself is not evaluated, only the campaign design.
This is the same gap the 27 July issue named and it has not moved.

\item \textbf{Ultrafine and speciated metrics: improving, from nothing.} 6 of 131 records
(4.6\%) use UFP or particle-number exposure and 3 use speciated composition. That is up
from near-zero, entirely because the aerosol-science journals are now being read. It
remains badly out of proportion to the deposition physics: the fraction of deposited
particle \emph{number} that reaches the alveolar region is dominated by exactly the size
range that 95\% of this week's records do not measure.

\item \textbf{Geographic coverage vs.\ attributable burden: the gap narrowed.} 13 of 131
records (10\%) come from South Asia or sub-Saharan Africa --- India 7, Ghana, Uganda,
South Africa, Bangladesh and a three-country Gambia/Kenya/Mozambique study. Against
China's 42 (32\%) this is still inverted relative to attributable burden, but it is the
best week on this measure so far, and two of the records (Nyoni, Ghana AIRCLARE) are
primary measurement rather than secondary analysis.

\item \textbf{Calibration transferability: reported, and it fails.} Three records tested
transfer explicitly. Ginsburg et al.\ (29 Jul) found US-derived PurpleAir corrections do
not transfer outside North America across 1{,}013 globally co-located sensors, with
region-specific corrections needed to meet EPA 24\,h targets. Kumpika et al.\ report
$R^2 = 0.974$ but on single-site, single-season, single-model data --- co-located fit, not
held-out deployment. Qian et al.\ report chamber performance only. \textbf{No record this
week reported held-out deployment performance at a site not used in fitting.}

\item \textbf{Negative and null results: present, which is a good sign.} Zhu et al.'s null
PM$_{2.5}$ pneumonia RR, Blanco et al.'s null UFP--cognition estimate, Yang et al.'s null
PM$\times$heat interaction and the AFRI-c filtration process evaluation all published as
nulls. A corpus with no nulls would be a reporting-bias signal; this one has four in six
days, which is consistent with a literature that is at least sometimes publishing what it
finds.
\end{itemize}

\vspace{3mm}
\begin{keybox}
{\sffamily\bfseries\footnotesize\color{Deep}Where this is heading --- judgement, not
measurement}\\[1.2mm]
\footnotesize
The methodological demand the field is converging on is \textbf{source- and
size-resolved exposure with a validated error term}. Three separate literatures arrived
at it independently this week: epidemiology (Dales's wildfire-vs-background split, Zhu's
null mass term), exposure modelling (Saeedi and Blanco on validation and design), and
aerosol science (Zhou L on number-dominated indoor emissions, Mishra on predicting
oxidative potential from composition). None of them cites the others. That convergence is
the strongest signal in the week and I would put moderate-to-high confidence on it
persisting.

What would change the picture: a well-powered cohort finding a robust total-mass
association after adjusting for source and composition would undercut the whole line, and
none of this week's records tested that. The mass-null results here come from
single-pollutant models over correlated combustion tracers, which is a weak basis for
demoting mass. \textbf{I would not yet write ``PM$_{2.5}$ mass is the wrong metric'' in a
thesis introduction} --- I would write that its measured association is increasingly
sensitive to source attribution and model specification, which is a defensible and much
narrower claim.

For a distributed low-cost PM sensing programme the operational reading is specific.
First, \textbf{number-based channels are now worth the extra cost}: the week produced two
independent results (Bakhoor UFP emission rates exceeding sidestream cigarette smoke;
6 UFP records against essentially zero the previous month) that a mass-only network
cannot see. Second, \textbf{validate on held-out sites, not held-out samples} --- Saeedi's
217\%$\to$79\% collapse is what random-CV validation costs, and it applies directly to
any network calibration reported as ``$R^2$ against co-located reference''. Third,
\textbf{instrument ageing and firmware behaviour are first-order error terms}: Kumpika's
duty-cycling result and Le's CL61 laser-power/calibration-factor finding both describe
degradation that a deployed node cannot self-detect. A co-location rotation schedule is
not optional infrastructure.

Confidence that this week was informative: high --- but note that the single largest
change in the corpus was a fix to \emph{this pipeline}, not to the literature. Nine of
the twenty sensing records exist in this rollup because a harvest leg was built on
Friday. The literature did not change; the instrument reading it did.
\end{keybox}

\clearpage

% ==================================================================== DIGEST
\section*{\sffamily\bfseries\color{Deep}Paper-level digest --- the week by subtopic}
\vspace{-1.5mm}{\color{Teal}\rule{\linewidth}{1.1pt}}\vspace{2mm}

\footnotesize
Every record of the week, grouped by subtopic and ordered by first author within group.
\textbf{123 of 131 records carry the full critical summary written when they were first
screened} --- design, effect size with interval, the limitation that threatens the
inference, and the line on relevance to distributed sensing. Those summaries are
reproduced \emph{verbatim} from the archived daily issues; nothing was re-written for this
rollup, and the index is rebuilt by parsing \texttt{issues/digest\_*.tex} rather than
retyped.

The remaining \textbf{8 records were surfaced by the Crossref by-journal leg on 1 August
and have no abstract} deposited in Crossref, Europe PMC or PubMed. They are listed on
title, journal and DOI alone. \textbf{No summary is offered for them and the absence of
criticism is not approval} --- they are unread, and are marked as such.

\vspace{3mm}

\input{digest_body}

\clearpage

% ==================================================================== TABLE
\section*{\sffamily\bfseries\color{Deep}Pooled corpus register}
\vspace{-1.5mm}{\color{Teal}\rule{\linewidth}{1.1pt}}\vspace{2.5mm}

{\fontsize{6.7}{8.4}\selectfont
\setlength{\tabcolsep}{2.4pt}
\begin{longtable}{@{}>{\RaggedRight}p{25mm}>{\RaggedRight}p{41mm}>{\RaggedRight}p{27mm}%
>{\RaggedRight}p{23mm}>{\RaggedRight}p{18mm}>{\RaggedRight}p{15mm}@{}}
\toprule
\textsf{\textbf{First author}} & \textsf{\textbf{Focus}} & \textsf{\textbf{Journal}} &
\textsf{\textbf{Design}} & \textsf{\textbf{Metric}} & \textsf{\textbf{Region}}\\
\midrule
\endfirsthead
\toprule
\textsf{\textbf{First author}} & \textsf{\textbf{Focus}} & \textsf{\textbf{Journal}} &
\textsf{\textbf{Design}} & \textsf{\textbf{Metric}} & \textsf{\textbf{Region}}\\
\midrule
\endhead
\multicolumn{6}{@{}l@{}}{\cellcolor{Mist}\textsf{\textbf{\textcolor{Deep}{Neuro / mental health}}}}\\[0.6mm]
\href{https://doi.org/10.1001/jamanetworkopen.2026.25015}{Dales} & Wildfire-sourced PM$_{2.5}$ and ED visits for migraine and other primary headache & \textit{JAMA Netw Open} & Case-crossover & PM$_{2.5}$ (wildfire) & Canada\\
\href{https://doi.org/10.1093/ageing/afag215}{Glover} & Umbrella meta-review of urban environmental exposures and cognitive health & \textit{Age Ageing} & Umbrella review / meta-analysis & PM$_{2.5}$ & Global\\
\href{https://doi.org/10.1016/j.nsa.2026.107019}{Goossens} & In utero air pollution exposure and pre-/postnatal brain development & \textit{Neurosci Appl} & Narrative review & PM$_{2.5}$ / PM$_{10}$ / PM$_{\mathrm{coarse}}$ & Global\\
\href{https://doi.org/10.1016/j.isci.2026.116617}{Ma} & PM$_{2.5}$-induced neurotoxicity and depression: exposure risks and mechanisms & \textit{iScience} & Narrative review & PM$_{2.5}$ + components & Global\\
\href{https://doi.org/10.3760/cma.j.cn112150-20250804-00752}{Pan} & Maternal PM exposure and offspring neurodevelopment: effects and mechanisms & \textit{Zhonghua Yu Fang Yi Xue Za Zhi} & Narrative review & PM (bulk) & China\\
\addlinespace[0.9mm]
\multicolumn{6}{@{}l@{}}{\cellcolor{Mist}\textsf{\textbf{\textcolor{Teal}{Cardiovascular \& metabolic}}}}\\[0.6mm]
\href{https://doi.org/10.3760/cma.j.cn112338-20251120-00843}{Huang} & PM$_{2.5}$ and dyslipidaemia in rural Chinese children; dietary effect modification & \textit{Zhonghua Liu Xing Bing Xue Za Zhi} & Cross-sectional (multistage) & PM$_{2.5}$ & China\\
\href{https://doi.org/10.32394/pe/218966}{Kazmierska} & Air pollution and cardiovascular disease: review of current data & \textit{Przegl Epidemiol} & Narrative review & PM$_{2.5}$ / PM$_{10}$ & Global\\
\href{https://doi.org/10.1007/s00125-026-06814-2}{Xi} & Air pollution, built environment and temperature vs. prediabetes phenotypes & \textit{Diabetologia} & Prospective cohort & PM$_{2.5}$ & Germany\\
\href{https://doi.org/10.1161/JAHA.125.046917}{Yang} & Short-term PM$_{2.5}$ and hospitalisation/ED use in atrial fibrillation & \textit{J Am Heart Assoc} & Case-crossover & PM$_{2.5}$ (wildfire-influenced) & USA\\
\addlinespace[0.9mm]
\multicolumn{6}{@{}l@{}}{\cellcolor{Mist}\textsf{\textbf{\textcolor{Sky}{Reproductive \& developmental}}}}\\[0.6mm]
\href{https://doi.org/10.1016/j.jes.2025.10.043}{Chen} & Prenatal air pollution, fetal growth trajectory phenotypes and child neurodevelopment & \textit{J Environ Sci (China)} & Prospective cohort & PM$_{2.5}$ / PM$_{10}$ & China\\
\href{https://doi.org/10.1016/j.ecoenv.2026.120566}{Huang} & Joint prenatal bisphenol + air-pollutant mixture exposure and preschool IQ & \textit{Ecotoxicol Environ Saf} & Prospective cohort (mixtures) & PM$_{2.5}$ / PM$_{10}$ & China\\
\href{https://doi.org/10.1007/s40572-026-00552-8}{Peterson} & Ambient air pollution, PAHs and pubertal development: critical review & \textit{Curr Environ Health Rep} & Critical review & PM$_{2.5}$ components & Global\\
\href{https://doi.org/10.1016/j.envres.2026.125313}{Xie} & Long-term PM$_{2.5}$ and timing of pubertal development in boys and girls & \textit{Environ Res} & Prospective cohort & PM$_{2.5}$ & China\\
\addlinespace[0.9mm]
\multicolumn{6}{@{}l@{}}{\cellcolor{Mist}\textsf{\textbf{\textcolor{Sage}{Respiratory \& allergic}}}}\\[0.6mm]
\href{https://doi.org/10.1007/s00484-026-03289-z}{Bayar Muluk} & Climate change and otolaryngology: airway disease impacts and emission reduction & \textit{Int J Biometeorol} & Narrative review & PM$_{2.5}$ & Global\\
\href{https://doi.org/10.4103/mgr.MEDGASRES-D-25-00164}{Liao} & Short-term air pollution, respiratory symptoms and lung function in asthma & \textit{Med Gas Res} & Panel study & PM$_{2.5}$ / PM$_{10}$ & China\\
\href{https://doi.org/10.3390/epidemiologia7040097}{Sanchez de Prada} & Ambient PM and asthma-coded admissions in children 0-3 y, Spain & \textit{Epidemiologia} & Ecological panel & PM$_{2.5}$ / PM$_{10}$ & Spain\\
\href{https://doi.org/10.1016/j.jhazmat.2026.143076}{Yang} & Allergenic vs. toxic risk of airborne seed hairs compared with ambient PM$_{2.5}$ & \textit{J Hazard Mater} & Field measurement & PM$_{2.5}$ + bioaerosol & China\\
\addlinespace[0.9mm]
\multicolumn{6}{@{}l@{}}{\cellcolor{Mist}\textsf{\textbf{\textcolor{Violet}{Mechanistic toxicology}}}}\\[0.6mm]
\href{https://doi.org/10.1016/j.jes.2025.12.074}{Ding} & PM$_{2.5}$ promotes vascular calcification via miR-27b-5p/HES1 signalling & \textit{J Environ Sci (China)} & In vitro + organoid + murine & PM$_{2.5}$ & China\\
\href{https://doi.org/10.1007/s00360-026-01698-5}{Maraschi} & Osmoregulatory disruption and metal bioaccumulation from settleable atmospheric PM & \textit{J Comp Physiol B} & Controlled exposure (ecotox) & Settleable APM & Brazil\\
\href{https://doi.org/10.1016/j.jes.2025.11.018}{Wu} & Macrophage NLRP3 inflammasome in real-ambient PM-induced hepatic glucose dysregulation & \textit{J Environ Sci (China)} & Real-ambient murine + co-culture & Real-ambient PM & China\\
\addlinespace[0.9mm]
\multicolumn{6}{@{}l@{}}{\cellcolor{Mist}\textsf{\textbf{\textcolor{Clay}{Occupational \& indoor}}}}\\[0.6mm]
\href{https://doi.org/10.3389/fmedt.2026.1796828}{Chen} & Nested suction instrument to cut intraoperative surgical-smoke PM$_{2.5}$ & \textit{Front Med Technol} & Bench / device experiment & PM$_{2.5}$ (surgical smoke) & China\\
\href{https://doi.org/10.1093/annweh/wxag055}{Cole} & Determinants of respirable crystalline silica exposure in Australian tunnelling & \textit{Ann Work Expo Health} & Observational exposure study & Respirable dust / RCS & Australia\\
\href{https://doi.org/10.1016/j.jhazmat.2026.143081}{Ding} & Stage-resolved hair Zn-isotope dynamics as a firefighting combustion-exposure tracer & \textit{J Hazard Mater} & Repeated-measures biomarker & Combustion-derived PM & China\\
\href{https://doi.org/10.3390/jox16040119}{Ramsden} & Xenobiotic hazards in aircraft cabin air & \textit{J Xenobiot} & Review + risk economics & Ultrafine particles & Global\\
\addlinespace[0.9mm]
\multicolumn{6}{@{}l@{}}{\cellcolor{Mist}\textsf{\textbf{\textcolor{Amber}{Sensing, forecasting \& instrumentation}}}}\\[0.6mm]
\href{https://doi.org/10.1016/j.jenvman.2026.130515}{Shi} & Adaptive spatiotemporal graph transformer for multi-site multi-step PM$_{2.5}$ forecasting & \textit{J Environ Manage} & Deep-learning model & PM$_{2.5}$ & China / India\\
\href{https://doi.org/10.3390/bios16070380}{Tsuchiya} & Seasonal and indoor-outdoor characteristics of fluorescent aerosol particles in offices & \textit{Biosensors} & Real-time sensor campaign & Bioaerosol / coarse & Japan\\
\addlinespace[0.9mm]
\multicolumn{6}{@{}l@{}}{\cellcolor{Mist}\textsf{\textbf{\textcolor{Amber}{Exposure assessment \& modelling}}}}\\[0.6mm]
\href{https://doi.org/10.1029/2025GH001621}{Solorzano-Araque} & Multiscale exposure modelling: LOTOS-EUROS + mobility surveys + particle toxicity & \textit{GeoHealth} & CTM + mobility integration & PM$_{2.5}$ / PM$_{10}$ & Colombia\\
\addlinespace[0.9mm]
\multicolumn{6}{@{}l@{}}{\cellcolor{Mist}\textsf{\textbf{\textcolor{Coral}{Burden, policy \& mitigation}}}}\\[0.6mm]
\href{https://doi.org/10.1002/hsr2.72886}{Karamzad} & Ischaemic heart disease burden in MENA 1990-2021; PM attributable fraction 28.2\% & \textit{Health Sci Rep} & GBD secondary analysis & Ambient PM & MENA\\
\href{https://doi.org/10.1371/journal.pone.0347989}{Rees} & HEPA filtration to reduce respiratory infection in care homes & \textit{PLoS One} & RCT process evaluation & Filtered PM & UK\\
\href{https://doi.org/10.1016/j.jes.2026.01.009}{Shan} & Policy-driven emission reductions and post-pandemic rebound, high-resolution inventories & \textit{J Environ Sci (China)} & Emission inventory & PM emissions & China\\
\addlinespace[0.9mm]
\multicolumn{6}{@{}l@{}}{\cellcolor{Mist}\textsf{\textbf{\textcolor{Slate}{Other clinical endpoints}}}}\\[0.6mm]
\href{https://doi.org/10.3390/dj14070440}{Lupita} & PM$_{2.5}$, preventive-care capacity and caries/edentulism burden across the EU & \textit{Dent J} & Ecological panel & PM$_{2.5}$ & EU-27\\
\href{https://doi.org/10.1016/j.jdin.2026.06.008}{Park} & Long-term PM exposure and risk of androgenetic alopecia & \textit{JAAD Int} & Nationwide cohort & PM$_{2.5}$ / PM$_{10}$ & South Korea\\
\href{https://doi.org/10.3171/2026.3.PEDS2615}{Sun} & Environmental exposure and area deprivation vs. paediatric hydrocephalus outcomes & \textit{J Neurosurg Pediatr} & Registry cohort + geospatial & PM$_{2.5}$ (satellite-derived) & USA\\
\addlinespace[0.9mm]
\bottomrule
\end{longtable}}


\vspace{4mm}
\begin{keybox}
{\sffamily\bfseries\footnotesize\color{Deep}Provenance and method}\\[1.2mm]
{\fontsize{7.2}{8.8}\selectfont
This rollup is an \textbf{aggregation only}. It pools \texttt{state/corpus/*.json},
\texttt{state/metrics.csv} and \texttt{state/trials.json} over
\PERIODSTART{}--\PERIODEND{} via \texttt{PMRW\_START}/\texttt{PMRW\_END} and
\textbf{re-queries no API}, so every record was screened and deduplicated in the daily
issue that first carried it. Per-record provenance, source-query windows, exclusions and
known metadata defects are documented in those daily issues and are not restated here.
\\[1.0mm]
\textbf{Daily issues pooled: 27, 28, 29, 30, 31 July and 1 August 2026 --- six
consecutive issues, no gap.}
\\[1.0mm]
\textbf{Why the window starts on 27 July and not 26 July.} The weekly cadence is
Sunday-to-Saturday, so the nominal window is 26 July--1 August. The 26 July slot in
\texttt{state/corpus/} does not hold a 26 July issue: it holds the \textbf{1--26 July
retrospective backfill} (164 records, entry window \texttt{2026-07-01 -> 2026-07-26})
harvested for the July monthly rollup. Pooling it here would put a month of records
inside a week, inflate every count by a factor of 2.3 and double-count material already
published in \texttt{Reports/monthly/PM-Research-Watch-Monthly\_2026-07-31.pdf}. It is
therefore excluded and the pooled window is stated as 27 July--1 August throughout.
The 164 backfill records are not lost --- they are in the July monthly in full.
\\[1.0mm]
Effect estimates are transcribed from abstracts and use heterogeneous exposure contrasts
--- see the forest-plot caption; \textbf{no meta-analytic summary is computed and no
$I^2$ is reported}, because the estimates are not exchangeable. Trend and inflection
statements are judgements over a six-day window and are not statistical tests. The
``where this is heading'' box is explicitly judgement and is labelled as such. Content
derived from PubMed, Europe PMC and Crossref metadata; abstracts remain the copyright of
their respective publishers.\par}
\end{keybox}

\end{document}
